% \iffalse meta-comment
%
% Copyright (C) 2017- by Yanshuo Chu <yanshuoc@gmail.com>
%
% This file may be distributed and/or modified under the
% conditions of the LaTeX Project Public License, either version 1.3a
% of this license or (at your option) any later version.
% The latest version of this license is in:
%
% http://www.latex-project.org/lppl.txt
%
% and version 1.3a or later is part of all distributions of LaTeX
% version 2004/10/01 or later.
%
% \fi
%
% ^^A ======================================================================
% ^^A hithesis 用户手册
% ^^A
% ^^A 只放面向用户的说明（模板介绍 / 安装 / 使用说明），不含任何实现代码，
% ^^A 也不参与 docstrip：hithesis.ins 从不读这个文件。
% ^^A 由 hithesis.dtx 的 driver 通过 \DocInput 引入，渲染成手册的前半部分。
% ^^A ======================================================================
%
% \GetFileInfo{\jobname.dtx}
%
%
% \def\indexname{索引}
% \def\glossaryname{修改记录}
% \IndexPrologue{\section{\indexname}}
% \GlossaryPrologue{\section{\glossaryname}}
%
% \title{\bfseries\color{violet}\hithesis：哈尔滨工业大学学位论文模板}
% \changes{v3.2a}{0000/00/00}{封面的作者信息加上 GitHub 账号，并注明当前主要维护者}
% \author{{\fangsong
%     \hitauthor{初砚硕}{\hitmail{yanshuoc@gmail.com}}{dustincys}，
%     \hitauthor{孙忠豪}{\hitmail{syvshc@foxmail.com}}{syvshc}，
%     \hitauthor{石佳乐}{\hitmail{mingzzx@outlook.com}}{MingZZX}，
%     \hitauthor{李云鹏}{\hitmail{ln3dqagt03gnqxkqt0@163.com}}{xiF616}，
%     \hitauthor{支卓}{\hitmail{xeonacid@hit.edu.cn}}{Xeonacid}，
%     \hitauthor{吴俊豪}{\hitmail{andrew.junhao.wu@gmail.com}}{SchrodingerBlume}}
%   \\[10pt]
%   \parbox{0.82\textwidth}{\centering\normalsize\normalfont
%     自 v3.2a 起，\hithesis 完成了基于 \pkg{expl3} 的底层重构，
%     目前维护工作主要由吴俊豪承担；使用中如遇问题，建议优先与他联系。}}
% \date{v\fileversion\ (\filedate)}
% \maketitle\thispagestyle{empty}
% \begin{abstract}\noindent
% \indent 本宏包为哈尔滨工业大学（一校三区）本科、硕士、博士研究生开题报告、
%   中期报告、中英文毕业论文（设计），以及博士后出站报告的模板。
% \end{abstract}
%
% \vskip2cm
% \changes{v3.2a}{0000/00/00}{更新免责声明}
% \def\abstractname{免责声明}
% \begin{abstract}
% \noindent
% \begin{enumerate}
% \item 本模板由作者根据\hit 教务处发布的\UGR 及\hit 研究生院发布的\PGR 编写，
%   供\hit 学生撰写毕业材料、博士后撰写出站材料时使用。
% \item \hit 教务处和研究生院仅提供毕业论文写作指南，一般不提供官方模板（包括
%   MS Word 模板），也未授权任何第三方模板为官方模板\footnote{目前深圳校区的
%   官方本科 \LaTeX 模板基于本宏包修改而来，但这并不代表本宏包获得了官方认可。}。
%   因此本模板仅是写作指南的一种参考实现，不保证评审专家不会提出修改意见，
%   因使用本模板而产生的任何问题，\hithesis 维护团队不承担任何责任。
% \item 本模板的发布遵守 \LaTeX\ Project Public License，使用前请认真阅读协议
%   内容。任何个人或组织以本模板为基础进行修改、扩展而生成的专用模板，均须严格
%   遵守 \LaTeX\ Project Public License 协议，因违反协议而引起的任何纠纷，
%   \hithesis 维护团队不承担任何责任。
% \end{enumerate}
% \end{abstract}
%
%
% \clearpage
% \pagestyle{fancy}
% \begin{multicols}{2}[
%   \setlength{\columnseprule}{.4pt}
%   \setlength{\columnsep}{18pt}]
%   \tableofcontents
% \end{multicols}
% \clearpage
%
% \section{模板介绍}
% \hithesis\ ({\color{violet}H}arbin {\color{violet}I}nstitute of {\color{violet}T}echnology \LaTeX\
% {\color{violet}Thesis} Template) 是为了帮助\hit 学生撰写开题报告、中期报告、
% 毕业论文（设计），以及博士后撰写出站报告而编写的 \LaTeX 文档模板。
%
% 本文档将尽量翔实地介绍模板的使用方法。如有不清楚之处，可以参考示例文档，或按
% 第~\ref{sec:howtoask}~节的说明提问。欢迎参与完善本手册或贡献代码；有意长期参与
% 维护的同学，请联系 \hitref{SchrodingerBlume}。
%
% \subsection*{文档记号的说明}
% 本文档采用不同字体表示不同内容。无衬线字体表示宏包名称，如
% \pkg{xeCJK} 宏包、\cls{hit-thesis} 文档类等；等宽字体表示代码或
% 文件名，如 \cs{hitsetup} 命令、\env{abstract} 环境、\TeX{} 文档
% \file{thesis.tex} 等；带有尖括号的楷体（或西文斜体）表示命令参数，
% 如 \option{模板选项}、\option{English title} 等。在使用时，参数两侧
% 的尖括号不必输入。示例代码进行了语法高亮处理，以方便阅读。
% 在用户手册中，带有\textcolor{violet}{紫色}侧边线的为 \LaTeX{} 代码，
% 而带有\textcolor{purple}{红色}侧边线的则为命令行代码，请注意区分。
%
% \note[注意：]{模板的作用在于减少论文写作过程中格式调整的时间。前提是遵守模板的
% 用法，否则即便用了 \hithesis\ 也难以保证输出的论文符合学校规范。}
%
%
% \section{安装}
% \label{sec:installation}
% ^^A 未来
% ^^A \hithesis\ 将已经包含在主要的 \TeX\ 发行版中，一般不需要安装，可以利用发行版自
% ^^A 带更新工具自动更新。阅读文档可以使用以下命令：
% ^^A \begin{shell}
% ^^A $ texdoc hithesis
% ^^A \end{shell} 
% \textbf{目前还没有将最新版的 \hithesis\ 上传到 CTAN，所以不要使用 \texttt{texdoc} 来打开文档。}
%
% 可以在相关 QQ 群内下载编译好的模板文件和文档，但是文档还没有更新完毕，有一些内容是过时的，有一些更新的功能没有添加到文档中，等写好文档了就传 CTAN，又挖了个坑。
%
% 如果要使用开发版，需自己下载，\hithesis\ 相关链接是 
% \begin{itemize}
% \item github：\href{https://github.com/hithesis/hithesis}
% {https://github.com/hithesis/hithesis}
% \item github dev 分支：\href{https://github.com/hithesis/hithesis/tree/dev}
% {https://github.com/hithesis/hithesis/tree/dev}
% \end{itemize}
% 其中 dev 分支中的为最新开发版，有可能有还没解决的 bug （目前来看好像不太可能），如果没有额外描述，请尽量使用 dev 分支的模板
% \subsection{模板的组成}
% 下表列出了 \hithesis\ 的主要文件及其功能介绍：
%
% \begin{longtable}{lp{8cm}}
% \toprule
% {\heiti 文件（夹）} & {\heiti 功能描述}\\\midrule
% \endfirsthead
% \midrule
% {\heiti 文件（夹）} & {\heiti 功能描述}\\\midrule
% \endhead
% \endfoot
% \endlastfoot
% \file{hithesis.ins} & \textsc{DocStrip} 驱动文件（开发用） \\
% \file{hithesis.dtx} & \textsc{DocStrip} 主源文件（开发用）\\
% \file{hit-*.dtx} & 各个类与共用模块的源文件（开发用）\\
% \file{hit-eps.dtx} & EPS 图像数据（校徽、封面、示例插图等，从主源剥离）\\
% \file{dtx-style.sty} & \textsc{DocStrip} 样式文件（开发用）\\\midrule
% \file{hit-thesis.cls} & 学位论文类文件\\
% \file{hit-report.cls} & 开题和中期报告类文件 \\
% \file{hithesisbook.cls} & 转发到 \file{hit-thesis.cls} 的空壳，v4.1a 移除 \\
% \file{hithesisart.cls} & 转发到 \file{hit-report.cls} 的空壳，v4.1a 移除 \\
% \file{hithesisartplus.cls} & 转发到 \file{hit-report.cls} 的空壳，v4.1a 移除 \\
% \file{hithesis.bst} & 参考文献样式文件\\
% \file{hitszthesis.bst} & 深圳校区参考文献样式文件 \\\midrule
% \file{hithesis.ist} & 索引样式文件\\\midrule
% \file{examples/} & 模板示例文件夹 \\
% \texttt{├─ }\file{hitart/} & 开题和中期报告文件夹\\
% \texttt{│\phantom{─ }├─ }\file{reports/} & 开题和中期报告文件夹 \\
% \texttt{│\phantom{─ }└─ }\file{reportplus/} & 深圳校区博士中期报告文件夹 \\
% \texttt{└─ }\file{hitbook/} & 毕业论文文件夹\\
% \texttt{\phantom{├─ }├─ }\file{chinese/} & 中文毕业论文文件夹 \\
% \texttt{\phantom{├─ }└─ }\file{english/} & 英文毕业论文文件夹 \\\midrule
% \file{Makefile} & Makefile\\
% \file{latexmkrc} & latexmk 配置文件 \\
% \file{README.md} & Readme\\
% \file{\textbf{hithesis.pdf}} & 用户手册（本文档）\\\bottomrule
% \end{longtable}
%
% 几点说明：
% \begin{itemize}
% \item 所有的 \file{hithesis*.cls} 以及 \file{hithesis.bst}、
%   \file{hitszthesis.bst}、\file{hithesis.ist} 都可由 \file{hithesis.ins} 和 \file{hithesis.dtx} 生成。
% \item 使用前阅读文档： \file{hithesis.pdf} 。
% \item 使用 \file{hithesis.ins} 生成样式文件后， \file{examples/} 文件夹中的 \file{chinese/}，\file{english/} ，\file{reports/} 与 \file{reportplus/} 文件夹均可以单独使用。
% \end{itemize}
%
% \subsection{代码组织}
% \label{sec:modules}
% \changes{v3.2a}{0000/00/00}{类文件内部按职责分区重写，用户接口不变}
% v3.2a 把类文件的实现按职责分成了若干区块（选项、字体、元信息、封面、章节标题、参考文献、
% 页眉页脚、目录等），每个区块只管一件事。这些区块在生成 \file{*.cls} 时按固定顺序拼接，
% 最终仍是三个独立的类文件，用户接口与此前完全一致。
%
% 分区只影响源码的组织方式，用户不会接触到。维护者若要了解各区块的职责划分与新增区块的
% 步骤，见仓库根目录的 \file{CONTRIBUTING.md}。
%
% \subsection{expl3 与空格}
% \label{sec:expl3-space}
% \changes{v3.2a}{0000/00/00}{手册补一节说明 expl3 语法下空格的处理，
%   \file{hit-thesis.dtx} 里早有一处 \cs{ref} 指向这里}
% v3.2a 起类文件从头到尾都在 \pkg{expl3} 的 catcode régime 里，由开头的
% \cs{ProvidesExplClass} 打开，中间不再有成对的 \cs{ExplSyntaxOn} 与
% \cs{ExplSyntaxOff}。这套 catcode 下空格的含义和平时不一样，改代码时有三件事要留神。
%
% \textbf{字面空格会被吞。} \texttt{text=\{150true mm\}} 照抄进去会变成
% \texttt{150truemm}，要空格得写 \texttt{\~{}}。控制词后面的空格是名字终结符，
% 本来就会被吸收，不用管。行末的换行同样被忽略，多行的文本值里原本靠换行产生的
% 那个空格要补 \texttt{\~{}}，原本靠空行产生的 \cs{par} 要显式写出来。
%
% \textbf{\texttt{\~{}} 的含义反过来了。} 平时它是活动字符，表示不断行空格；
% \pkg{expl3} 下它是普通空格。原有代码里的 \texttt{\~{}} 照抄进来，不断行空格
% 就退化成了普通空格，要保持原意得写 \cs{nobreakspace}\{\}。
%
% \textbf{控制空格不在此列}，它在两种语法下是同一个控制
% 序列，照写即可。
%
% 另有两处例外，必须退回 2e 的 catcode：走裸 \cs{input} 的加载器（如
% \cs{usetikzlibrary}），以及 \pkg{etoolbox} 的 \cs{patchcmd} 一族——它们会把
% 文件或宏体按当前 catcode 重扫，字面空格照样被吞。v3.2a 已经把这两类都消掉了，
% 类文件里一处 \cs{ExplSyntaxOn} 都没有。
%
% \subsection{生成模板}
% \label{sec:generate-cls}
% \note[提示：]{若使用 \TeX 发行版自带的 \hithesis\ 或 Github/OSChina
% 上发布的版本，可忽略此节，直接阅读第 \ref{sec:generate-thesis} 节。若下载
% CTAN 包或者 Github/OSChina 开发代码，请阅读本节了解生成模板文件的步骤。}
%
% 模板解压缩后生成文件夹 \file{hithesis-vX.Y.Z} 或 \file{hithesis-vX.YZ}\footnote{\texttt{vX.Y.Z} 与 \texttt{vX.YZ} 为版本号。}，
% 其中包括：模板源文件（\file{hithesis.ins} 和 \file{hithesis.dtx}），参考文献
% 样式 \file{hithesis.bst}，示例文档
% （\file{thesis.tex}\footnote{自 v3.2a 起，原 \file{hithesis.sty} 的代码已按职责并入
% 类文件内部。为了让导言区还写着 \cs{usepackage}\marg{hithesis} 的老文档编得过，
% 会生成一个空壳 \file{hithesis-compat.sty}，并把 \file{hithesis.sty} 这个名字重定向
% 过去——用户升级时常常只覆盖 \file{.cls}，旧宏包留在目录里会先被找到。空壳本身
% 不做任何事，v4.1 一并移除。}，
% \file{front/}、\file{body/}、
% \file{back/} 和 \file{figures/}）。在使用之前需要先生成模
% 板文件和配置文件（具体命令细节请参考 \file{README.md} 和 \file{Makefile}）：
%
% \begin{shell}
% $ cd hithesis-vX.Y.Z
% # 生成类文件与配置文件，并分发到各示例目录
% $ make cls
%
% # 下面的命令用来生成用户手册，可以不执行
% $ make doc
% \end{shell}
%
% 没有 \texttt{make} 的话，可以直接调 \TeX{}：
%
% \begin{shell}
% $ latex hithesis.ins
% \end{shell}
%
% 这条命令只把生成物放在当前目录。示例目录是独立可用的，需要把用到的那份
% 类文件、\file{.bst} 与 \file{.eps} 复制过去，例如：
%
% \begin{shell}
% $ cp hit-thesis.cls hithesis.ist hithesis.bst \
%      hitszthesis.bst *.eps examples/hitbook/chinese/
% $ cp golfer.eps examples/hitbook/chinese/figures/
% \end{shell}
%
% \texttt{make cls} 做的就是这件事，所以有 \texttt{make} 时不必手工复制。
%
% \subsection{生成论文}
% \label{sec:generate-thesis}
% 本节介绍几种常见的生成论文的方法。用户可根据自己的情况选择。
%
% \subsubsection{\XeLaTeX}
% \label{sec:xelatex}
% 很多用户对 \LaTeX\ 命令执行的次数不太清楚。一个基本的原则是多次运行 \LaTeX\ 命
% 令直至不再出现警告。下面给出生成示例文档的详细过程（\texttt{\#} 开头的行为注
% 释），首先来看推荐的 \texttt{xelatex} 方式：
% \begin{shell}
% # 1. 发现里面的引用关系，文件后缀 .tex 可以省略
% $ xelatex thesis
%
% # 2. 编译参考文件源文件，生成 bbl 文件
% $ bibtex thesis
%
% # 3. 下面解决引用
% $ xelatex thesis
% $ xelatex thesis   # 如果不需要生成索引此时生成完整的 pdf 文件
% $ splitindex thesis -- -s hithesis.ist  # 自动生成索引
% $ xelatex thesis.tex
% \end{shell}
%
% \subsubsection{latexmk}
% \label{sec:latexmk}
% \texttt{latexmk} 命令支持全自动生成 \LaTeX\ 编写的文档，并且支持使用不同的工具
% 链来进行生成，它会自动运行多次工具直到交叉引用都被解决。下面给出了一个用
% \texttt{latexmk} 调用 \texttt{xelatex} 生成最终文档的示例：
% \begin{shell}
% # 一句话就够了！
% $ latexmk -xelatex thesis
% \end{shell}
%
% \subsubsection{make}
% \label{sec:make}
% \note[提示：]{若要使用 \texttt{make} 编译，需自行下载模板。因为 \TeX\ 发行版中
% 的 \file{Makefile} 不在当前目录。}
%
% 上面的方法虽然不复杂，但是每次都输入还是非常罗嗦，所以 \hithesis\ 提供了一
% 个 \file{Makefile}：
%
% 仓库根目录下的 \file{Makefile} 负责生成类文件与用户手册：
% \begin{shell}
% $ make clean
% $ make cls       # 生成 hit-thesis.cls 等
% $ make doc       # 生成说明文档 hithesis.pdf
% \end{shell}
% 各示例文件夹（如 \file{examples/hitbook/chinese/}）下另有一份 \file{Makefile}，用于编译示例文档：
% \begin{shell}
% $ make thesis    # hitbook 示例生成 thesis.pdf；hitart 示例对应目标为 make report
% \end{shell}
%
% \hithesis\ 的 \file{Makefile} 默认用 \texttt{latexmk} 调用\texttt{xelatex} 编
% 译，此外还支持直接用 \texttt{xelatex} 编译。如有需要可修
% 改 \file{Makefile} 开头的参数或通过命令行传递参数（请参看 \file{README.md}），
% 进一步还可以修改 \file{latexmkrc} 进行定制。
%
% \subsection{升级}
% \label{sec:updgrade}
% \hithesis\ 升级非常简单，可以通过 \TeX 发行版的包管理工具自动更新发行版，也可
% 以下载最新的开发版，
% 将 \file{hithesis.ins}，\file{hithesis.dtx} 拷贝至工作目
% 录覆盖相应的文件，然后运行：
% \begin{shell}
% $ latex hithesis.ins
% \end{shell}
%
% 生成新的类文件即可。也可以直接拷
% 贝 \file{hit-thesis.cls}
% 与 \file{hithesis.ist}，免去上面命令的执行。
%
%
% \subsection{从 v3.1e 升上来要注意的}
% \label{sec:breaking}
% \changes{v3.2a}{0000/00/00}{手册加一节列出相对 v3.1e 的破坏性变更}
% v3.2a 把各个类的底层全部迁到 \pkg{expl3}，对外接口大体不变，但有几处确实变了。
% 下面按“会不会编不过”排序。
%
% \subsubsection{编不过的}
%
% \begin{itemize}
% \item \textbf{删掉了 \cs{dachu}、\cs{banxiaosi}、\cs{dawu} 三个字号。}
%   这三个不在通行的中文字号里，手册的字号对照表从来没列过。用到的改成相邻的
%   标准字号：\cs{dachu}（58bp）改 \cs{chuhao}，\cs{banxiaosi}（13bp）改
%   \cs{sihao} 或 \cs{xiaosi}，\cs{dawu}（11bp）改 \cs{xiaosi} 或 \cs{wuhao}。
% \item \textbf{删掉了 \option{pifootnote} 选项与 \pkg{pifont}。}
%   带圈脚注号现在统一由 \cs{quan} 排，先取当前宋体自带的带圈字符，超出字体可用
%   范围才画圈。要 Dingbats 那套 ➀➁➂ 的自己 \cs{usepackage}\marg{pifont}。
% \item \textbf{不再预载 \pkg{tikz} 的 \texttt{calc} 库。}
%   自己的插图里写了 \texttt{(\$…\$)} 坐标运算的，在导言区补一句
%   \cs{usetikzlibrary}\marg{calc}。

% \item \textbf{最低 \TeX~Live 版本提到 2022。}
% \end{itemize}
%
% \subsubsection{能编过但要留意的}
%
% \begin{itemize}
% \changes{v3.2a}{0000/00/00}{手册记一条类改名的升级说明}
% \item \textbf{文档类改名了。}
%   类文件与源文件同名：\file{hit-thesis.cls} 出学位论文，\file{hit-report.cls}
%   出开题与中期报告。三个旧名各留一个空壳，把选项原样转过去：
%
%   \begin{longtable}{ll}
%   \toprule
%   旧写法 & 改成 \\\midrule
%   \endfirsthead
%   \toprule
%   旧写法 & 改成 \\\midrule
%   \endhead
%   \cs{documentclass}\marg{hithesisbook} & \cs{documentclass}\marg{hit-thesis} \\
%   \cs{documentclass}\marg{hithesisart} & \cs{documentclass}\marg{hit-report} \\
%   \cs{documentclass}\marg{hithesisartplus} &
%     \cs{documentclass}\oarg{type=doctor,campus=shenzhen,stage=midterm}\marg{hit-report} \\\bottomrule
%   \end{longtable}
%
%   深圳博士中期原先自成一个文档类 \file{hithesisartplus.cls}，与
%   \file{hithesisart.cls} 九成代码逐字相同，现在是报告类里
%   \texttt{campus=shenzhen,type=doctor,stage=midterm} 这一格，与其余十七格并列。
%   空壳会补齐这三个选项，所以不带选项的旧写法也照编。
%
%   三个空壳都在 v4.1a（2027 年主要版本）移除，届时不改的文档会直接编不过。
% \changes{v3.2a}{0000/00/00}{手册记一条目录条目与书签去掉字间距的说明}
% \item \textbf{目录条目与 PDF 书签里的“结论”不再带撑开的空隙。}
%   规范要求“结论”“致谢”“目录”这类两字标题撑开排，原先撑开写在固定词汇的
%   取值里，同一个值又供目录条目与书签用，书签里读出来带着空隙。现在取值只留
%   文字，撑开由排标题的地方做。变的是目录条目、页眉与书签。
%
%   撑多少由 \option{const} 族的 \option{heading-spread} 定，默认一个全角字宽
%   （\texttt{1em}）——Word 版模板里那两个字之间是两个半角空格，宋体的半角空格
%   是半个字宽。原先写的是 \cs{ccwd}，那是“字宽加字距”，开着 \option{zijv}
%   时比一个字宽多 5.6\%，所以三个标题的字间距各窄了 1bp。博士后那一档规范
%   不撑开，取空值。
% \changes{v3.2a}{0000/00/00}{手册记一条博士后封面字段标签对齐的说明}
% \item \textbf{博士后封面的四个字段标签现在真的对齐了。}
%   “博士后姓名”“合作导师”“流动站名称”“专业名称”本该是同一列宽，原先靠手写
%   字间距凑。凑的时候字数按字宽算、间距按 \cs{ccwd} 算，而开了 \option{zijv}
%   之后 \cs{ccwd} 是字宽的 1.0558 倍，两条式子把这个差额放大成不同的量，
%   四个标签一个 99.94bp、一个 100.72bp，值那一列对不齐。现在改成分散对齐到
%   7 个字宽，四行都是 98.00bp，比原先窄 1.5bp。
% \changes{v3.2a}{0000/00/00}{手册记一条摘要标题一并归拢的说明}
% \item \textbf{“摘要”标题与“结论”“致谢”“目录”同一条路了。}
%   它原先也是把撑开写在取值里（\cs{cabstractcname}），同样漏进目录条目与书签。
%   现在一并走 \option{heading-spread}，字间距一个全角字宽，目录条目与书签取原文。
% \changes{v3.2a}{0000/00/00}{手册记一条双语题注短标题与 \cs{label} 的说明}
% \item \textbf{双语题注的 \cs{caption}\oarg{短标题} 与 \cs{label} 都能用了。}
%   \cs{caption}\oarg{短标题}\marg{中文题}\oarg{英文题} 里的短标题原先没进插图索引，
%   而是被悄悄当成了 \cs{label}——\pkg{ccaption} 的 \cs{bicaption} 可选参数就是标签，
%   模板接错了位置。现在短标题进索引，与只写中文题时一致；\cs{label} 照普通
%   \LaTeX{} 的写法放在 \cs{caption} 之后。
%
%   靠 \cs{caption}\oarg{名字} 当标签用的文档要改成 \cs{label}\marg{名字}，
%   否则那条引用会变成 \texttt{??}。
% \changes{v3.2a}{0000/00/00}{手册记一条几何符号归中文字体的说明}
% \item \textbf{中文里的 \texttt{□}、\texttt{○}、\texttt{△} 排得出来了。}
%   \pkg{xeCJK} 按 Unicode 区块分中西文，“几何图形”那一块归了西文，所以中文正文里
%   写 \texttt{□}（规范原文用它占位，如“攻读□士学位”）会去找西文字体。
%   \pkg{TeX~Gyre~Termes} 里没有这个字，结果不是排成豆腐块，是整个字没了，
%   只在日志里留一句 \texttt{Missing character}。现在这六个符号归中文字体，正常排出。
% \item \textbf{\cs{usepackage}\marg{hithesis} 可以删了。}
%   宏包的代码已按职责并入类文件。为了让老文档编得过，\file{hithesis.sty} 这个名字
%   被重定向到一个空壳，它什么都不做，只发一条提示；目录里留着旧版宏包也不会撞车。
%   空壳在 v4.1 一并移除，届时不删这一行的文档会直接编不过。
% \item \textbf{\cs{bicaption} 不再需要。}
%   英文题注改成 \cs{caption} 的尾随可选参数，写
%   \cs{caption}\marg{中文题}\oarg{英文题}。\cs{bicaption} 的三种旧写法（\pkg{ccaption}
%   的五参数原形、本模板的三参数与四参数）都还能用，各发一条建议改写的警告。
% \item \textbf{子图从 \pkg{subfigure} 换成了 \pkg{subcaption}。}
%   \cs{subfigure}、\cs{subtable}、\cs{subref} 三个旧命令仍然能用，
%   \cs{ifsubcaphang} 等五个开关会自动转成对应的 \cs{captionsetup}，但都会发一条
%   建议改写的警告。
% \item \textbf{定理从 \pkg{ntheorem} 换成了 \pkg{thmtools}。}
%   \cs{theorembodyfont}、\cs{theoremheaderfont}、\cs{theoremsymbol}、
%   \cs{theoremseparator} 会转成 \cs{hitsetup} 里的对应键；\cs{theoremprework}、
%   \cs{theorempostwork}、\cs{theoremclass}、\cs{listtheorems}、\cs{thref} 没有
%   对应实现，会提示改用 \pkg{thmtools} 的写法。
% \item \textbf{类选项从 \pkg{kvoptions} 换成了 l3keys。}
%   排版类的选项写在 \cs{documentclass} 的可选参数里仍然生效，但会提示改用
%   \cs{hitsetup}，这条支持将在 v4.1 移除。
% \item \textbf{\env{publist} 改用 \pkg{enumitem} 实现。}
%   用法与版面都没变，仍然是 \cs{item} 逐条列。
% \item \textbf{\option{citetwo}\texttt{=comma} 现在真的生效了。}
%   这个选项从 2022 年加进来就一直是坏的（补丁打不上但吞掉了失败），修好之后
%   两个相邻连续序号会排成 \texttt{[1,2]} 而不是 \texttt{[1-2]}。原先误以为它
%   没作用的，升级后版面会变。
% \item \textbf{\cs{hitsetup}\marg{bib=\{adjacent-cite=…\}} 与
%   \cs{hitsetup}\marg{style=\{ziju=…\}} 原先会报错}，现在修好了。
% \item \textbf{中文的斜体改成倾斜，不再是楷体。}
%   \changes{v3.2a}{0000/00/00}{手册记一条中文斜体的改动}
%   \cs{textit} 与 \cs{emph} 下的中文原先显示为楷体，现在是倾斜的正文字体。
%   要楷体请直接写 \cs{kaishu}。只影响写了 \option{font-zh} 的文档，不写这个选项
%   的中文字体仍由 \CTeX{} 设置，行为不变。
% \item \textbf{\option{font} 与 \option{cjk-font} 改名 \option{font-en} 与
%   \option{font-zh}。}
%   \changes{v3.2a}{0000/00/00}{手册记一条 \option{font} 与 \option{cjk-font} 的改名}
%   与模板其余选项的中英后缀一致。旧名仍然生效，只在日志里提示一次，将在 v4.1 移除。
% \item \textbf{\file{hitszthesis.bst} 补上了 \texttt{@masterthesis} 的容错。}
%   \changes{v3.2a}{0000/00/00}{手册记一条 \file{hitszthesis.bst} 的
%     \texttt{@masterthesis} 容错}
%   \BibTeX{} 的标准条目类型是 \texttt{@mastersthesis}，写漏 \texttt{s} 的
%   \file{hithesis.bst} 一直认，\file{hitszthesis.bst} 原先不认，那条文献会不声不响
%   地从文献表里消失。现在两个样式一致。升级后深圳校区的文献表可能多出几条。
% \item \textbf{art 与 artplus 的脚注改成了与 book 一致的样式。}
%   \changes{v3.2a}{0000/00/00}{art 与 artplus 的脚注改成与 book 一致：带圈编号、
%     宋体小五、悬挂 1.5 字符}
%   编号由上标阿拉伯数字改成带圈数字，正文改成宋体小五，悬挂缩进 1.5 字符，分割线
%   宽 0.3\cs{textwidth}。我工规范对各个类的脚注要求本来就是一样的，原先只有
%   book 照做。
% \end{itemize}
%
% \subsubsection{新增的}
%
% \begin{itemize}
% \item \option{font-en} 与 \option{font-zh} 两个选项，西文与中文字体可以分别指定，
%   见第 \ref{sec:chinese-fonts} 节。
% \item \option{math-font} 与 \option{math-style} 两个选项，数学字体走
%   \pkg{unicode-math}，见第 \ref{sec:math-fonts} 节。
% \item \env{eqdenote} 环境，排公式下面的“式中……”注释表，见第
%   \ref{sec:eqdenote} 节。
% \item \option{style/glue} 选项 art 与 artplus 也有了，与 book 一致。
% \item \env{publist} 列表 art 与 artplus 也有了，与 book 一致。
% \end{itemize}
%
% \section{使用说明}
% \label{sec:usage}
% 本手册假定用户已经能处理一般的 \LaTeX\ 文档，并对 \BibTeX\ 有一定了解。如果
% 从来没有接触过 \TeX\ 和 \LaTeX，建议先学习相关的基础知识。
%
% \subsection{关于提问}
% \label{sec:howtoask}
% 按照优先级推荐提问的位置如下（三个 QQ 群不分先后）：
%
% \begin{itemize}
% \item \href{http://github.com/hithesis/hithesis/issues}{Github Issues}
% ^^A \item \href{https://git.oschina.net/dustincys/hithesis/issues}{OSChina Issues}
% \item hithesis讨论区：259959600 （人满）
% \item 窝工山hithesis派：851792460 （人满）
% \item HiThesis 分盟：704864357
% \end{itemize}
%
% \subsection{示例文件}
% \label{sec:userguide}
% 模板核心文件有两
% 项：\file{hit-thesis.cls} 与 \file{hithesis.bst}，但是如果
% 没有示例文档用户会发现很难下手。所以推荐新用户从模板自带的示例文档入手，里面包
% 括了论文写作用到的所有命令及其使用方法，只需要用自己的内容进行相应替换就可以。
% 对于不清楚的命令可以查阅本手册。下面的例子描述了模板中章节的组织形式，来自于示
% 例文档，具体内容可以参考模板附带的 \file{thesis.tex}。
%
% \lstinputlisting[style=lstStyleLaTeX]{examples/hitbook/english/thesis.tex}
% \lstinputlisting[style=lstStyleLaTeX]{examples/hitbook/chinese/thesis.tex}
% \lstinputlisting[style=lstStyleLaTeX]{examples/hitart/reports/report.tex}
% \lstinputlisting[style=lstStyleLaTeX]{examples/hitart/reportplus/report.tex}
%
% \subsection{选哪个文档类}
% \label{sec:class}
% \changes{v3.2a}{0000/00/00}{手册补上派发类 \file{hithesis.cls} 与 \option{kind}}
%
% 有两个真正的文档类：\file{hit-thesis.cls} 排学位论文，\file{hit-report.cls}
% 排开题与中期报告。直接写它们一直有效。
%
% \DescribeMacro{kind}
% 另外有一个派发类 \file{hithesis.cls}，\option{kind} 选交付物种类，它按这个
% 转发给上面两个之一：
%
% \begin{latex}
% \documentclass[kind=thesis,  type=doctor, campus=harbin]{hithesis}
% \documentclass[kind=opening, type=doctor, campus=harbin]{hithesis}
% \documentclass[kind=midterm, type=doctor, campus=harbin]{hithesis}
% \end{latex}
%
% \begin{longtable}{lll}
% \toprule
% \option{kind} & 转发给 & 等价于 \\\midrule
% \endfirsthead
% \toprule
% \option{kind} & 转发给 & 等价于 \\\midrule
% \endhead
% \texttt{thesis}（默认） & \file{hit-thesis} & \cs{documentclass}\marg{hit-thesis} \\
% \texttt{opening} & \file{hit-report} &
%   \cs{documentclass}\oarg{stage=opening}\marg{hit-report} \\
% \texttt{midterm} & \file{hit-report} &
%   \cs{documentclass}\oarg{stage=midterm}\marg{hit-report} \\\bottomrule
% \end{longtable}
%
% 好处是一个学生三年里的开题、中期、学位论文可以放在同一个目录里，元信息、
% 参考文献库、图片目录写一次，三份主文件只有第一行不同。开题报告要填的字段
% 有九个，其中六个跟学位论文重名。
%
% \option{kind} 只能写在 \cs{documentclass} 上，挪不进导言区：\cs{LoadClass}
% 得在类文件里执行完，而 \cs{hitsetup} 是导言区的事，那时候基础类早装上了。
% 所以做不到“一个主文件靠一行开关切三种交付物”，能做到的是一个目录里三个
% 主文件，每个只有第一行不同。
%
% \option{kind} 写错了会有提示，并按 \texttt{thesis} 兜底，不硬失败——
% 装类失败的话后面几百个错都指不到病根。
%
% \subsection{论文选项}
% \label{sec:option}
%
% 论文选项，就是在 \file{thesis.tex} 文件的开头，非注释的第一行的方括号中填写的选项，示例见上节。
% 这几项决定版面，必须在加载基础类之前定下来，所以只能写在 \cs{documentclass} 里，
% 其余配置一律走 \cs{hitsetup}。
%
% \changes{v3.2a}{0000/00/00}{\option{language} 改名 \option{lang}，取值改用
%   \texttt{zh} 与 \texttt{en}，与字段后缀一致}
% \begin{longtable}{lll}
% \toprule
% 选项 & 取值 & 管什么 \\\midrule
% \endfirsthead
% \toprule
% 选项 & 取值 & 管什么 \\\midrule
% \endhead
% \option{type}    & \texttt{bachelor} / \texttt{master} / \texttt{doctor} /
%                    \texttt{postdoc} & 学位 \\
% \option{campus}  & \texttt{harbin} / \texttt{shenzhen} / \texttt{weihai} & 校区 \\
% \option{stage}   & \texttt{opening} / \texttt{midterm} & 报告类的阶段 \\
% \option{kind}    & \texttt{thesis} / \texttt{opening} / \texttt{midterm} &
%                    交付物种类，只有派发类 \file{hithesis} 认，见 \ref{sec:class} 节 \\
% \option{lang}    & \texttt{zh} / \texttt{en} & 正文语言，默认 \texttt{zh} \\
% \option{fontset} & 见 \CTeX\ 手册 & 字体集 \\\bottomrule
% \end{longtable}
%
% \option{lang} 只管摘要之后的正文。封面、扉页、原创性声明这些前置页面无论哪种
% 语言都照中文格式排，英文封面本来就是中文论文也有的那一页。
%
% \subsection{配置项的组织}
% \label{sec:keyfamily}
% \changes{v3.2a}{0000/00/00}{\cs{hitsetup} 的配置项按职责分成六个族，旧的扁平
%   写法进入一年兼容期}
% \cs{hitsetup} 的配置项按职责分成六个族：
%
% \begin{longtable}{ll}
% \toprule
% 族 & 管什么 \\\midrule
% \endfirsthead
% \toprule
% 族 & 管什么 \\\midrule
% \endhead
% \option{info}  & 封面与声明用到的元信息，如标题、作者、导师 \\
% \option{style} & 排版开关，如章标题加粗、页码带横线 \\
% \option{toc}   & 目录，如是否排英文目录、是否插空白页 \\
% \option{bib}   & 参考文献，如著者大小写、连续引用的分隔符 \\
% \option{const} & 固定词汇，见第~\ref{sec:const}~节 \\
% \option{theorem} & 定理环境的字体与间距，见第~\ref{sec:math}~节 \\\bottomrule
% \end{longtable}
%
% 写法是嵌套的键值表：
% \begin{latex}
% \hitsetup{
%   info = {
%     zh = { title = {论文标题}, author = {张三} },
%   },
%   toc = { english = true },
%   style = { chapter-bold = false },
% }
% \end{latex}
%
% 族名之下的键一律用连字符分词，如 \option{chapter-bold}、\option{adjacent-cite}。
%
% \subsubsection{三态取值}
% \label{sec:tristate}
% 有些配置项的合适取值随学位而变，这类配置项取三种值：
%
% \begin{longtable}{ll}
% \toprule
% 取值 & 含义 \\\midrule
% \endfirsthead
% \toprule
% 取值 & 含义 \\\midrule
% \endhead
% \texttt{auto} & 由模板按学位和校区决定，默认值 \\
% \texttt{true} / \texttt{false} & 一律开或一律关 \\
% 分学位映射 & 逐个学位指定，没列到的学位仍按 \texttt{auto} \\\bottomrule
% \end{longtable}
%
% \begin{latex}
% \hitsetup{
%   toc = { english = auto },                      % 模板决定
%   toc = { blank = false },                       % 一律不插
%   toc = { four = { master = true, doctor = true } },  % 只对硕博开
% }
% \end{latex}
%
% 取三态的有 \option{toc/english}、\option{toc/blank}、\option{toc/four}、
% \option{style/headrule}、\option{style/pagenumberline/mainmatter} 与
% \option{style/pagenumberline/frontmatter}。
%
% \subsubsection{旧写法}
% \label{sec:deprecated}
% v3.2a 之前所有配置项都写在同一层，且有一批名字带 \texttt{bs} 前缀或不带连字符。
% 旧写法仍然可用，编译时会在日志里提示新写法，兼容期一年，
% v4.1（2027~年主要版本）移除。对照表：
%
% \begin{longtable}{ll}
% \toprule
% 旧名 & 新写法 \\\midrule
% \endfirsthead
% \toprule
% 旧名 & 新写法 \\\midrule
% \endhead
% \option{arialtoc} & \option{toc} = \{ \option{sans-font} \} \\
% \option{tocfour} & \option{toc} = \{ \option{four} \} \\
% \option{tocblank} & \option{toc} = \{ \option{blank} \} \\
% \option{engtoc} & \option{toc} = \{ \option{english} \} \\
% \option{arialtitle} & \option{style} = \{ \option{sans-title} \} \\
% \option{chapterbold} & \option{style} = \{ \option{chapter-bold} \} \\
% \option{chapterhang} & \option{style} = \{ \option{chapter-hang} \} \\
% \option{capcenterlast} & \option{style} = \{ \option{caption-center-last} \} \\
% \option{subcapcenterlast} & \option{style} = \{ \option{subcaption-center-last} \} \\
% \option{absupper} & \option{style} = \{ \option{abstract-upper} \} \\
% \option{etitleauto} & \option{style} = \{ \option{etitle-auto} \} \\
% \option{glue} & \option{style} = \{ \option{glue} \} \\
% \option{openright} & \option{style} = \{ \option{open-right} \} \\
% \option{library} & \option{style} = \{ \option{library} \} \\
% \option{raggedbottom} & \option{style} = \{ \option{ragged-bottom} \} \\
% \option{zijv} & \option{style} = \{ \option{ziju} \} \\
% \option{bsheadrule} & \option{style} = \{ \option{headrule} \} \\
% \option{bsmainpagenumberline} & \option{style} = \{ \option{pagenumberline} = \{ \option{mainmatter} \} \} \\
% \option{bsfrontpagenumberline} & \option{style} = \{ \option{pagenumberline} = \{ \option{frontmatter} \} \} \\
% \option{citetwo} & \option{bib} = \{ \option{adjacent-cite} \} \\
% \option{splitbibitem} & \option{bib} = \{ \option{split-item} \} \\
% \option{biblabelalign} & \option{bib} = \{ \option{label-align} \} \\
% \option{bibauthorupper} & \option{bib} = \{ \option{author-case} \} \\
% \option{bibmaxauthor} & \option{bib} = \{ \option{max-author} \} \\\bottomrule
% \end{longtable}
%
% 元信息字段在 v3.2a 也改成了英文名。原先的 \texttt{c}\ldots{} 与 \texttt{e}\ldots{}
% 前缀取消，中英文两版改用 \texttt{-zh} 与 \texttt{-en} 后缀，与 \option{font-zh}、
% \option{font-en} 一致。旧名同样保留到 v4.1。
%
% \begin{longtable}{ll}
% \toprule
% 旧名 & 新写法 \\\midrule
% \endfirsthead
% \toprule
% 旧名 & 新写法 \\\midrule
% \endhead
% \option{caffil} & \option{info} = \{ \option{affiliation-zh} \} \\
% \option{eaffil} & \option{info} = \{ \option{affiliation-en} \} \\
% \option{cassosupervisor} & \option{info} = \{ \option{associate-supervisor-zh} \} \\
% \option{eassosupervisor} & \option{info} = \{ \option{associate-supervisor-en} \} \\
% \option{cauthor} & \option{info} = \{ \option{author-zh} \} \\
% \option{eauthor} & \option{info} = \{ \option{author-en} \} \\
% \option{cclassid} & \option{info} = \{ \option{class-id} \} \\
% \option{intclassifiedindex} & \option{info} = \{ \option{classified-index-international} \} \\
% \option{natclassifiedindex} & \option{info} = \{ \option{classified-index-national} \} \\
% \option{ccosupervisor} & \option{info} = \{ \option{co-supervisor-zh} \} \\
% \option{ecosupervisor} & \option{info} = \{ \option{co-supervisor-en} \} \\
% \option{firstpagecdate} & \option{info} = \{ \option{cover-date-zh} \} \\
% \option{ctitlecover} & \option{info} = \{ \option{cover-title-zh} \} \\
% \option{cdate} & \option{info} = \{ \option{date-zh} \} \\
% \option{edate} & \option{info} = \{ \option{date-en} \} \\
% \option{cxueke} & \option{info} = \{ \option{discipline-zh} \} \\
% \option{exueke} & \option{info} = \{ \option{discipline-en} \} \\
% \option{cenrolldate} & \option{info} = \{ \option{enroll-date-zh} \} \\
% \option{ckeywords} & \option{info} = \{ \option{keywords-zh} \} \\
% \option{ekeywords} & \option{info} = \{ \option{keywords-en} \} \\
% \option{cnumber} & \option{info} = \{ \option{report-number} \} \\
% \option{csubmitdate} & \option{info} = \{ \option{report-submit-date-zh} \} \\
% \option{ctype} & \option{info} = \{ \option{report-type-zh} \} \\
% \option{cenddate} & \option{info} = \{ \option{research-end-date-zh} \} \\
% \option{cstartdate} & \option{info} = \{ \option{research-start-date-zh} \} \\
% \option{statesecrets} & \option{info} = \{ \option{secret-level} \} \\
% \option{csubject} & \option{info} = \{ \option{specialty-zh} \} \\
% \option{esubject} & \option{info} = \{ \option{specialty-en} \} \\
% \option{cpositionname} & \option{info} = \{ \option{station-zh} \} \\
% \option{cstudentid} & \option{info} = \{ \option{student-id} \} \\
% \option{cstudenttype} & \option{info} = \{ \option{student-type-zh} \} \\
% \option{estudenttype} & \option{info} = \{ \option{student-type-en} \} \\
% \option{csubtitle} & \option{info} = \{ \option{subtitle-zh} \} \\
% \option{esubtitle} & \option{info} = \{ \option{subtitle-en} \} \\
% \option{csupervisor} & \option{info} = \{ \option{supervisor-zh} \} \\
% \option{esupervisor} & \option{info} = \{ \option{supervisor-en} \} \\
% \cs{bicaption}\marg{中文题}\marg{英文题} & \cs{caption}\marg{中文题}\oarg{英文题} \\
% \env{cabstract} & \env{abstract-zh} \\
% \env{eabstract} & \env{abstract-en} \\
% \option{ctitle} & \option{info} = \{ \option{title-zh} \} \\
% \option{ctitlecover} & \option{info} = \{ \option{title-zh} \} \\
% \option{csubtitle} & \option{title-zh} 的 \oarg{副标题} \\
% \option{esubtitle} & \option{title-en} 的 \oarg{副标题} \\
% \option{etitle} & \option{info} = \{ \option{title-en} \} \\
% \option{ctitleone} & \option{title-zh} 的第一个 \marg{分组} \\
% \option{ctitletwo} & \option{title-zh} 的第二个 \marg{分组} \\
% \option{cfinishdate} & \option{info} = \{ \option{work-finish-date-zh} \} \\\bottomrule
% \end{longtable}
%
% 排版类的配置项此前也可以写在 \cs{documentclass} 的可选参数里，同样进入兼容期，
% 请改用 \cs{hitsetup}。\cs{documentclass} 里保留 \option{type}、\option{campus}、
% \option{fontset}、\option{lang} 与 \option{stage}，它们决定版面，必须在加载
% 基础类之前就定下来。
%
% \subsection{西文字体}
% \changes{v3.2a}{0000/00/00}{西文字体独立成 \option{font-en} 选项}
% 西文字体由类选项 \option{font-en} 决定：
%
% \begin{longtable}{ll}
% \toprule
% 取值 & 字体 \\\midrule
% \endfirsthead
% \toprule
% 取值 & 字体 \\\midrule
% \endhead
% \texttt{auto} & 默认。见下 \\
% \texttt{times} & Times New Roman + Arial + Courier New \\
% \texttt{termes} & TeX Gyre Termes + Heros + Cursor \\
% \texttt{none} & 不设置，交给用户自己写 \\\bottomrule
% \end{longtable}
%
% \changes{v3.2a}{0000/00/00}{\option{font} 的 \texttt{auto} 一档改成除开源字体集外
%   一律优先 Times New Roman}
% \texttt{auto} 时按 \option{fontset} 决定：写的是开源字体集
% （\texttt{fandol}、\texttt{ubuntu}）就用 TeX Gyre Termes（Times 的开源克隆），
% 其余一律优先 Times New Roman——学校要求的就是它。机器上找不到会自动退回
% TeX Gyre，不会报错，所以 Linux 上装没装 Times New Roman 都能编。
%
% 之所以给开源字体集配开源西文，是因为写 \texttt{fandol} 或 \texttt{ubuntu} 的人
% 多半图的是“到哪台机器都编出一样的东西”，配一个装没装看运气的字体不合这个意思。
% 要在开源字体集上也用 Times New Roman，显式写 \option{font-en=times} 即可。
%
% \subsection{中文字体}
% \label{sec:chinese-fonts}
% 正确配置中文字体是使用模板的第一步。模板调用 \CTeX\ 宏包，只提供基于
% \pkg{xeCJK} 包，使用 \XeLaTeX\ 编译的方式。
% 关于如何使用字体命令、字号等等，属于模板格式范畴，在实现细节中讨论。
% 关于中文字体安装、配置的所有问题不在本模板讨论 范围。
%
% \changes{v3.2a}{0000/00/00}{中文字体独立成 \option{font-zh} 选项}
% 默认情况下中文字体由 \CTeX\ 按 \option{fontset} 配置，模板不插手。写了类选项
% \option{font-zh} 就改由模板自己设置，此时 \option{fontset} 只用来决定
% \option{font-zh} 取 \texttt{auto} 时选哪一档，不再转给 \CTeX。取值：
%
% \begin{longtable}{ll}
% \toprule
% 取值 & 字体 \\\midrule
% \endfirsthead
% \toprule
% 取值 & 字体 \\\midrule
% \endhead
% \texttt{auto} & 按 \option{fontset} 决定，都没写则按平台猜 \\
% \texttt{windows} & 中易宋体、黑体、楷体、仿宋 \\
% \texttt{mac} & 华文宋体、黑体、楷体、仿宋 \\
% \texttt{ubuntu} & 思源（Noto CJK 版），名字跟 \CTeX\ 的同名 fontset 对齐 \\
% \texttt{fandol} & Fandol，\TeX{} 发行版自带 \\
% \texttt{adobe} & Adobe Song/Hei/Kai/Fangsong Std \\
% \texttt{founder} & 方正书宋、黑体、楷体、仿宋 \\
% \texttt{hanyi} & 汉仪书宋二、中黑、楷体、仿宋 \\
% \texttt{none} & 不设置，用户自己 \cs{setCJKfamilyfont} 定义
%   \texttt{zhsong}、\texttt{zhhei}、\texttt{zhfs}、\texttt{zhkai} 四族 \\\bottomrule
% \end{longtable}
%
% \changes{v3.2a}{0000/00/00}{去掉模板自带的 siyuan fontset}
% 模板早年提供过一个 \texttt{siyuan} fontset（\file{ctex-fontset-siyuan.def}），
% v3.2a 起去掉了。思源与 Noto CJK 是同一套字，用 \option{font-zh=ubuntu} 或
% \pkg{ctex} 的 \option{fontset=ubuntu} 即可。旧文档里写的
% \option{fontset=siyuan} 会被 \pkg{ctex} 报成找不到字体集，然后回退到平台默认
% 字体继续编译。
%
% \subsection{相邻文献序号的连接符}
% \changes{v3.2a}{0000/00/00}{\option{adjacent-cite} 换实现，从 2022 年加进来
%   就没生效过}
% \option{bib} 族的 \option{adjacent-cite} 管两个相邻且连续的文献序号之间用什么
% 连接：
%
% \begin{longtable}{lll}
% \toprule
% 取值 & \cs{cite}\{a,b\}（序号 1、2） & \cs{cite}\{a,b,c\}（序号 1、2、3） \\\midrule
% \endfirsthead
% \toprule
% 取值 & \cs{cite}\{a,b\} & \cs{cite}\{a,b,c\} \\\midrule
% \endhead
% \texttt{endash} & \texttt{[1-2]}，默认 & \texttt{[1-3]} \\
% \texttt{comma} & \texttt{[1,2]} & \texttt{[1-3]} \\\bottomrule
% \end{longtable}
%
% 三个及以上连续的照样压成区间，只有正好两个的那种情况会换成逗号。不连续的序号
% 本来就是逗号分隔，不受这个选项影响。
%
% \begin{latex}
% \hitsetup{bib={adjacent-cite=comma}}
% \end{latex}
%
% \subsection{不再预载的 \pkg{tikz} 库}
% \changes{v3.2a}{0000/00/00}{不再预载 \pkg{tikz} 的 \texttt{calc} 库}
% 模板过去会替用户 \cs{usetikzlibrary}\marg{calc}。v3.2a 起不再预载：类文件
% 自己已经不用它了（封面那个边框改用 \pkg{tikz} 核心的
% \texttt{[shift=\{…\}]} 写法），而预载它会逼着类文件在这一句上退出 expl3
% 语法（\cs{usetikzlibrary} 走裸 \cs{input}，不恢复 catcode）。
%
% 自己的插图里用到 \texttt{(\$…\$)} 坐标运算的，在导言区补一句
% \cs{usetikzlibrary}\marg{calc} 即可。
%
% \subsection{数学字体}
% \label{sec:math-fonts}
% \changes{v3.2a}{0000/00/00}{数学引擎换成 \pkg{unicode-math}}
% 从 v3.2a 起数学公式走 \pkg{unicode-math}，数学字体由类选项 \option{math-font}
% 决定：
%
% \begin{longtable}{ll}
% \toprule
% 取值 & 字体 \\\midrule
% \endfirsthead
% \toprule
% 取值 & 字体 \\\midrule
% \endhead
% \texttt{auto} & 默认，等于 \texttt{xits} \\
% \texttt{xits} & XITS Math，Times 系，与正文字体最搭 \\
% \texttt{stix} & STIX Two Math \\
% \texttt{termes} & TeX Gyre Termes Math \\
% \texttt{newcm} & New Computer Modern Math \\
% \texttt{lm} & Latin Modern Math \\
% \texttt{libertinus} & Libertinus Math \\
% \texttt{newtx} & 老路子，8~位的 \pkg{newtxmath}，不走 \pkg{unicode-math} \\
% \texttt{none} & 不设置，用 \LaTeX{} 默认的 CM \\\bottomrule
% \end{longtable}
%
% 原来的 \option{newtxmath} 开关还在，默认值从 \texttt{true} 改成了
% \texttt{false}，写 \option{newtxmath=true} 等于 \option{math-font=newtx}。
%
% \option{math-style} 管数学排版风格：
%
% \begin{longtable}{ll}
% \toprule
% 取值 & 含义 \\\midrule
% \endfirsthead
% \toprule
% 取值 & 含义 \\\midrule
% \endhead
% \texttt{GB} & 默认。ISO 的基础上把 \cs{partial} 改直立 \\
% \texttt{ISO} & 按 ISO 80000-2，希腊字母一律倾斜 \\
% \texttt{TeX} & \TeX{} 传统写法，大写希腊字母直立 \\\bottomrule
% \end{longtable}
%
% 换引擎会影响已有文档的几处写法：
%
% \begin{itemize}
%   \item \cs{bm} 与 \cs{boldsymbol} 都接到了 \pkg{unicode-math} 的 \cs{symbf}
%     上。\pkg{unicode-math} 的符号是扩展 mathchar，\pkg{bm} 处理不了。
%   \item \cs{mathbf} 保持粗正体（接的是 \cs{symbfup}），与老版本一致。
%   \item \cs{mathrm}、\cs{mathit} 等仍从数学字体取字，而不是
%     \pkg{unicode-math} 默认的正文字体，这样 \cs{mathbf} 套希腊字母不会掉字。
%   \item \pkg{amssymb} 只在 \option{math-font=newtx} 与 \texttt{none} 下加载，
%     其余情况它提供的符号由数学字体自带。
% \end{itemize}
%
% \subsection{公式的物理量注释}
% \label{sec:eqdenote}
% \changes{v3.2a}{0000/00/00}{新增 \env{eqdenote} 环境}
% \DescribeEnv{eqdenote}
% 规范要求公式中第一次出现的物理量代号给出注释，注释的转行与破折号后第一个字
% 对齐。\env{eqdenote} 把这套排版包好了，每条注释写一个 \cs{item}：
%
% \begin{latex}
% \begin{equation}\label{eq:1}
%   a = b + c
% \end{equation}
% \begin{eqdenote}
%   \item $a$ & 追踪飞行器与目标飞行器之间的相对位置矢量；
%   \item $b$ & 推力所产生的加速度。
% \end{eqdenote}
% \end{latex}
%
% 符号与说明之间用 \texttt{\&} 隔开，符号按原样写，要排数学就自己加
% \texttt{\$}。末条不用写 \cs{\textbackslash\textbackslash}。
%
% 可选参数换引导词，传空则不排引导词但保留等宽的空位：
%
% \begin{latex}
% \begin{eqdenote}[其中] ... \end{eqdenote}
% \begin{eqdenote}[]     ... \end{eqdenote}
% \end{latex}
%
% 表格底下是 \pkg{xltabular}（\pkg{longtable} 加 \texttt{X} 列），放不下会自动
% 转页接排，不用手动拆。要手动拆时把后半段另起一个 \env{eqdenote}\texttt{[]}，
% 两段的左缘仍然对齐。
%
% 引导词与破折号可以改：
%
% \begin{latex}
% \hitsetup{const={zh={eqdenote-label={式中}}, en={eqdenote-label={where}},
%                  eqdenote-dash={——}}}
% \end{latex}
%
% 中文文档用 \texttt{eqdenote-label-zh}，\option{lang=en} 时用
% \texttt{eqdenote-label-en}。
%
% \subsection{前文}
% \label{sec:titlepage}
% 前文内容是正文之前，含封面、摘要、目录、符号表。
% 封面信息提供两种配置方法：一是通过统一设置命 令 \cs{hitsetup}
% 通过\emph{key=value} 形式完成；二是每个信息利用命令独立设置， 其中命令的名字跟
% \emph{key} 相同。两种方式可以交叉使用，并按顺序执行（即后来的设置会覆
% 盖前面的）。以 \texttt{c} 开头的命令跟中文相关，\texttt{e}
% 开头则为对应的英文。
%
% \DescribeMacro{\hitsetup}
% \cs{hitsetup} 用法与常见 \emph{key=value} 命令相同，如下：
% \begin{latex}
% \hitsetup{
%   key1 = value1,
%   key2 = {a value, with comma},
% }
% % 可以多次调用
% \hitsetup{
%   key3 = value3,
%   key1 = value11, % 覆盖 value1
% }
% \end{latex}
%
% \note[注意：]{配置项之间不能有空行，否则会报错。}
%
% 大多数命令的使用方法都是： \cs{command}\marg{arg}，例外者将具体指出。这些命令都
% 在示例文档的 \file{front/cover.tex} 中。
%
% \subsubsection{中英文两版的写法}
% \label{sec:setup-language}
% \changes{v3.2a}{0000/00/00}{手册改以 \option{zh} \option{en} 分组为主写法，
%   带后缀的名字降为等价写法}
% 中英文成对的字段写在 \option{zh} 与 \option{en} 两个分组里，字段名不带后缀：
%
% \begin{latex}
% \hitsetup{ info = {
%   zh = { title = {局部多孔质气体静压轴承关键技术的研究}, author = {于冬梅} },
%   en = { title = {Research on Key Technology of Aerostatic Bearing}, author = {Yu Dongmei} },
%   student-id = {1180800000},
% } }
% \end{latex}
%
% 这是推荐写法，示例文档里用的也是它。好处是一眼看得出“中文这一套”“英文这一套”，
% 加一种语言不必给每个字段都改名。不分语言的字段（\option{student-id}、
% \option{secret-level} 等）写在分组外面。
%
% \option{info} 与 \option{const} 两族都认这两个分组。
%
% 分组只是把字段名接上后缀再转发，所以带后缀的写法一样有效：
%
% \begin{latex}
% \hitsetup{ info = { title-zh = {...}, title-en = {...} } }
% \end{latex}
%
% 两者完全等价，混着写也行。带后缀的名字是键的本名，报错与改名提示里出现的是它，
% 下文各节也按本名列出，在分组里写的时候把后缀去掉就是了。
%
% \changes{v3.2a}{0000/00/00}{手册补上斜杠路径写法}
% 还有第三种写法，用斜杠把族名和字段名连起来，只设一两个字段时省一层括号：
%
% \begin{latex}
% \hitsetup{ info/zh/title = {...}, info/title-en = {...}, const/zh/school = {...} }
% \end{latex}
%
% 斜杠两边有没有空格都行。三种写法落到同一个键上，随便挑。
%
% 写错的分组会有提示：把不分语言的字段写进 \option{zh} 里，或者把一个根本不存在的
% 字段名写进去，都会说清是哪个键。斜杠写法同样。
%
% \subsubsection{密级}
% \label{sec:setup-secret}
% \DescribeMacro{secret-level}
% \DescribeMacro{classified-index-national}
% \DescribeMacro{classified-index-international}
% 定义秘密级别和国内国际索引号。
% \begin{latex}
% \hitsetup{
% secret-level={公开},
% classified-index-national={TM301.2},
% classified-index-international={62-5},
% }
% \end{latex}
%
% \subsubsection{论文标题}
% \myentry{论文标题}
% \DescribeMacro{title-zh}
% \DescribeMacro{title-en}
% \DescribeMacro{cover-title-zh}
% \DescribeMacro{title-line-one-zh}
% \DescribeMacro{title-line-two-zh}
% \DescribeMacro{subtitle-zh}
% \DescribeMacro{subtitle-en}
% 中英文标题，一个键管所有出现的地方。取值带两套断点：花括号分组给「题目栏是两条
% 固定宽度的线」那种版式分行（本科第二封面、博后第一封面、深圳博士中期封面），一组
% 一条线；\cs{\textbackslash\textbackslash} 给自由排的封面分行；方括号是副标题，
% 写了就排，不必再单独打开 \option{subtitle}。原创性声明的《\ldots》与 PDF 书签
% 取的是去掉断行的一整行。
%
% 副标题只出现在自由排的那几个封面上：硕博的第一、第二封面，哈尔滨本科第二封面，
% 以及英文封面。题目栏是两条固定线的三种封面——博后第一封面、威海本科第二封面、
% 深圳博士中期封面——版式里没有副标题这一栏，写了也不会排；报告类整个没有副标题。
% 这是学校版式定的，不是模板漏了。
% \begin{latex}
% \hitsetup{
%   zh={
%     cover-title={封面中文题目可断行},
%     title={论文中文题目},
%     subtitle={论文中文副题目（如果有）},
%     title-line-one={本科生论文中文题目上部分},
%     title-line-two={本科生论文中文题目下部分},
%   },
%   en={
%     title={Thesis English Title},
%     subtitle={Thesis English Sub-Title (if necessary)},
%   },
% }
% \end{latex}
%
% \subsubsection{作者姓名}
% \myentry{作者姓名}
% \DescribeMacro{author-zh}
% \DescribeMacro{author-en}
% 作者姓名。
% \begin{latex}
% \hitsetup{
%   zh={author={中文姓名}},
%   en={author={Name in Pinyin}}
% }
% \end{latex}
%
% \subsubsection{申请学位名称}
% \label{sec:degree}
% \myentry{学科名称}
% \DescribeMacro{discipline-zh}
% \DescribeMacro{discipline-en}
% 按照入学的培养计划中学科自行填写，具体学科名称不是本文档范畴。
% 学位名称（博士/硕士/学士）由文档类的 \option{type} 选项自动确定，无须手动设置。
%
% \begin{latex}
% \hitsetup{
%   zh={discipline={工学}},
%   en={discipline={Engineering}},
% }
% \end{latex}
%
% \subsubsection{院系名称}
% \myentry{院系名称}
% \DescribeMacro{affiliation-zh}
% \DescribeMacro{affiliation-en}
% 院系名称，同上，按照入学的培养计划中学科自行填写，具体院系名称不是本文档范畴。
% \begin{latex}
% \hitsetup{
%  zh={affiliation={机电工程学院}},
%  en={affiliation={\emultiline[t]{School of Mechatronics Engineering \\ Mechatronics Engineering}}},
% }
% \end{latex}
% \note[注意：]{个别学院英文名过长，使用以上方法自行换行。}
%
% \subsubsection{专业名称}
% \myentry{专业名称}
% \DescribeMacro{specialty-zh}
% \DescribeMacro{specialty-en}
% 专业名称，同上，按照入学的培养计划中学科自行填写，具体名称不是本文档范畴。
% \begin{latex}
% \hitsetup{
%  zh={specialty={机械制造及其自动化}},
%  en={specialty={Computer Science and Technology}},
% }
% \end{latex}
%
% \subsubsection{导师}
% \myentry{导师}
% \DescribeMacro{supervisor-zh}
% \DescribeMacro{supervisor-en}
% 直接导师。
% \begin{latex}
% \hitsetup{
%   zh={supervisor={导师~教授}},
%   en={supervisor={Supervisor}}
% }
% \end{latex}
%
% \myentry{副导师}
% \DescribeMacro{associate-supervisor-zh}
% \DescribeMacro{associate-supervisor-en}
% 副指导教师。
% \begin{latex}
% \hitsetup{
%   zh={associate-supervisor={副导师~副教授}},
%   en={associate-supervisor={2nd Boss}}
% }
% \end{latex}
%
% \myentry{联合导师}
% \DescribeMacro{co-supervisor-zh}
% \DescribeMacro{co-supervisor-en}
% 硕士生联合指导教师，博士生联合导师。
% \begin{latex}
% \hitsetup{
%   zh={co-supervisor={联合导师~教授}},
%   en={co-supervisor={3rd Boss}}
% }
% \end{latex}
%
% \subsubsection{成文日期}
% \myentry{日期}
% \DescribeMacro{date-zh}
% \DescribeMacro{date-en}
% 默认为当前时间，也可以自己指定。
% \begin{latex}
% \hitsetup{
%   zh={date={中文日期}},
%   en={date={English Date}},
% }
% \end{latex}
%
% \myentry{博士后日期}
% \DescribeMacro{research-start-date-zh}
% \DescribeMacro{research-end-date-zh}
% \DescribeMacro{report-submit-date-zh}
% \DescribeMacro{work-finish-date-zh}
% 博士后出站报告封面需要填写进站、出站、提交、完成等多个日期，对应命令分别为
% \cs{research-start-date-zh}、\cs{research-end-date-zh}、\cs{report-submit-date-zh}、\cs{work-finish-date-zh}。
% 其余学生类型的封面无须填写。
%
% \subsubsection{学生类型}
% \myentry{学生类型}
% \DescribeMacro{student-type-zh}
% \DescribeMacro{student-type-en}
% 非全日制教育申请学位者
%（同等学力人员）、（工程硕士）、（工商管理硕士）、
%（高级管理人员工商管理硕士）、（公共管理硕士）、（中职教师）、（高校教师）等，
% 具体要求按照入学的培养计划中学科自行填写，具体名称不是本文档范畴。
% \begin{latex}
% \hitsetup{
%   zh={student-type={同等学力人员}},
%   en={student-type={Master of Art}},
% }
% \end{latex}
%
% \subsubsection{学号}
% \myentry{学号}
% \DescribeMacro{student-id}
% 学号，具体要求按照入学的培养计划中学科自行填写。
% \begin{latex}
% \hitsetup{
%   student-id={9527},
% }
% \end{latex}
%
% \subsubsection{摘要}
% \myentry{摘要正文}
% \DescribeEnv{abstract-zh}
% \DescribeEnv{abstract-en}
% \note[说明：]{摘要正文只能用环境命令的形式，不支持 \cs{hitsetup}。}
%
% \begin{latex}
% \begin{abstract-zh}
%  摘要请写在这里...
% \end{abstract-zh}
%
% \begin{abstract-en}
%  Here comes the abstract in English...
% \end{abstract-en}
% \end{latex}
%
% \myentry{关键词}
% \DescribeMacro{keywords-zh}
% \DescribeMacro{keywords-en}
% 关键词用英文逗号分割写入相应的命令中，模板会解析各关键词并生成符合不同论文格式
% 要求的关键词格式。
% \begin{latex}
% \hitsetup{
%   zh={keywords={关键词 1, 关键词 2}},
%   en={keywords={keyword 1, keyword 2}}
% }
% \end{latex}
%
% \subsubsection{符号对照表}
% \DescribeEnv{denotation}
% 主要符号表环境，单独在文件 \file{front/denotation.tex} 中。
% 跟据\PGR\
% 示例中要求，我工符号表是table环境，示例文件如下，由于我工要求博士论文图表标题是双语，所以任何对单个标题的全局格式调整都会影响到双语标题，所以这里使用\cs{vspace}，具体见实现细节中的描述。
% \begin{latex}
% \begin{denotation}
% \begin{table}[h]%此处最好是h
% \caption{国际单位制中具有专门名称的导出单位}
% \vspace{0.5em}\centering\wuhao
% \begin{tabular}{ccccc}
% \toprule
% 量的名称&单位名称&单位符号&其它表示实例\\
% \midrule
% 频率&赫[兹]&Hz&s-1\\
% \bottomrule
% \end{tabular}
% \end{table}
% \end{denotation}
% \end{latex}
%
% \subsubsection{目录}
% 目录不需要用户干预，自动生成，具体命令已经写在 \file{thesis.tex} 中。
%
% \subsection{正文}
%
% \subsubsection{图和表}
% \hit\ 博士毕业论文要求使用中英双语图题、表题，这增加了维护难度。
% 因为现有唯一的方法是在已有的图题或表题的基础上再添加一行英语图题或表题。
% 两个题之间的距离具体多少不在\PGR\ 中要求。目前的方法是用户手动调节该距离。
% 关于图题\PGR\ 和\UGR\ 只规定了居中，并没有规定居中对齐。然而评审老师很多喜欢居
% 中且居中对齐。模板默认选项是居中但不居中对齐，如果需要居中对齐，那么需要在
% \file{thesis.tex}的文档类选项中设置选项capcenterlast=true。详细方法见前文的介绍
% 。
% \begin{hitrgu}
% 每个图均应有图题（由图序和图名组成），图题不宜有标点符号，图名在图序之后空1个
% 半角字符排写。图序按章编排，如第1章第一个插图的图号为“图1-1”。图题置于图下，硕
% 士论文只用中文，博士论文用中、英两种文字，居中书写，中文在上，要求中文用宋体5
% 号字，英文用Times New Roman 5号字。有图注或其它说明时应置于图题之上。引用图应
% 注明出处，在图题右上角加引用文献号。图中若有分图时，分图题置于分图之下或图题之
% 下，可以只用中文书写，分图号用a)、b)等表示。图中各部分说明应采用中文（引用的外
% 文图除外）或数字符号，各项文字说明置于图题之上（有分图时，置于分图题之上）。图
% 中文字用宋体、Times New Roman字体，字号尽量采用5号字（当字数较多时可用小5号字
% ，以清晰表达为原则，但在一个插图内字号要统一）。同一图内使用文字应统一。图表中
% 物理量、符号用斜体。
% \end{hitrgu}
% \DescribeMacro{\caption}
% 双语题注写成 \cs{caption}\marg{中文题}\oarg{英文题}：
%
% \begin{latex}
% \begin{figure}
%   \includegraphics{golfer}
%   \caption{打高尔夫球的人}[The person playing golf]
% \end{figure}
% \end{latex}
%
% 不给尾随的 \oarg{英文题} 就只排中文题。目录用的短标题仍是 \cs{caption} 原本那个
% 前置可选参数，两者位置不同不会混：\cs{caption}\oarg{短标题}\marg{中文题}\oarg{英文题}。
%
% 英文图表名（\texttt{Fig.}、\texttt{Table}）不用每次手写，取自 \option{const} 族的
% \option{figure-prefix-en} 与 \option{table-prefix-en}。
%
% \DescribeMacro{\hitsetup}
% \option{style/bilingual-caption} 是一键关掉所有双语题注的开关，默认开。关掉之后
% 英文题一概不排，同一份稿子投不同用途时不必逐条删。开着而某条题注没给英文题，
% 那一条按单语排，不报错也不留空。
%
% 中文题太长、插图索引里放不下时，用 \cs{caption} 原本那个前置可选参数放短标题：
% \begin{latex}
% \caption[短标题]{很长的中文题注……}[The person playing golf]
% \end{latex}
%
% \changes{v3.2a}{0000/00/00}{英文图表名收进 \option{const} 族，调用处不再写}
% 英文图表名取 \option{const} 族的 \option{figure-prefix-en} 与 \option{table-prefix-en}
% （默认 \texttt{Fig.} 与 \texttt{Table}，都已经带好负细空格，用来抵掉
% \cs{fnum@figure} 插的那个空格）。要换成别的名字：
% \begin{latex}
% \hitsetup{const={figure-prefix-en={Figure~}}}
% \end{latex}
%
% \DescribeMacro{\bicaption}
% \changes{v3.2a}{0000/00/00}{\cs{bicaption} 降为兼容写法，v4.1 移除}
% \cs{bicaption} 是 v3.2a 之前的写法，仍然可用但每种形式各提示一次，v4.1 移除。
% 三种形式都接得住：\pkg{ccaption} 原本的五参数
% \cs{bicaption}\oarg{标签}\marg{索引短标题}\marg{中文题}\marg{英文图表名}\marg{英文题}，
% 以及本模板之前的三参数与四参数。新写的文档请用 \cs{caption}。
%
% 单双语图题的方法如下，注释中说明。
% \begin{latex}
% \begin{figure}[htpb]
% \centering
% \includegraphics[width = 0.4\textwidth]{golfer}
% \caption[golfer1]{注意图中文字尽量用五号字
% }[The person playing golf]
% \end{figure}
% \end{latex}
% 单张单图题的格式如下，
% \begin{latex}
% \begin{figure}[h]
% \centering
% \includegraphics[width = 0.4\textwidth]{golfer}
% \caption{注意图中文字字号尽量用五号字}
% \end{figure}
% \end{latex}
% 并排图例。
% \begin{latex}
% \begin{figure}[htbp]
% \centering
% \begin{minipage}{0.4\textwidth}
% \centering
% \includegraphics[width=\textwidth]{golfer}
% \caption[golfer2]{打高尔夫球的人}[The person playing golf]
% \end{minipage}
% \begin{minipage}{0.4\textwidth}
% \centering
% \includegraphics[width=\textwidth]{golfer}
% \caption[golfer3]{打高尔夫球的人}[The person playing golf]
% \end{minipage}
% \end{figure}
% \end{latex}
% 子图图例。子图题只要一种语言时直接写在可选参数里：
% \begin{latex}
% \begin{figure}[htbp]
% \centering
% \subfigure[打高尔夫球的人~1]{\label{golfer41}\includegraphics[width=0.4\textwidth]{golfer}}
% \subfigure[打高尔夫球的人~2]{\label{golfer42}\includegraphics[width=0.4\textwidth]{golfer}}
% \caption[golfer4]{打高尔夫球的人}[The person playing golf]
% \end{figure}
% \end{latex}
%
% \DescribeMacro{\bisubcaption}
% 要中英文两行子图题，把 \cs{bisubcaption}\marg{中文}\marg{英文} 写进可选参数：
% \begin{latex}
% \begin{figure}[htbp]
% \centering
% \subfigure[\bisubcaption{打高尔夫球的人~1}{The person playing golf}]{%
%   \label{golfer41}\includegraphics[width=0.4\textwidth]{golfer}}
% \subfigure[\bisubcaption{打高尔夫球的人~2}{The person playing golf}]{%
%   \label{golfer42}\includegraphics[width=0.4\textwidth]{golfer}}
% \caption[golfer4]{打高尔夫球的人}[The person playing golf]
% \end{figure}
% \end{latex}
%
% \changes{v3.2a}{0000/00/00}{子图从 \pkg{subfigure} 换成 \pkg{subcaption}，
%   旧写法进入兼容期（\issue[284]）}
% v3.2a 之前子图由 \pkg{subfigure} 实现，双语子图题只能靠嵌套两层 \cs{subfigure}
% 再拿 \cs{addtocounter} 把计数器凑回来，标签也得用一个空的 \cs{subfigure} 单独占：
% \begin{latex}
% \subfigure{\label{golfer41}}\addtocounter{subfigure}{-2}
% \subfigure[The person playing golf]{\subfigure[打高尔夫球的人~1]{\includegraphics{golfer}}}
% \end{latex}
% \pkg{subfigure} 二十多年没人维护，与 \pkg{cleveref}、\pkg{hyperref} 配合有问题，
% 现已换成 \pkg{subcaption}。上面这种旧写法仍然可用，模板会自动转成新结构并在日志里
% 提示，v4.1（2027~年主要版本）移除。排出来的版面与旧版一致。
%
% \pkg{subfigure} 专有的开关与 \cs{listsubcaptions} 已经失效，用到会在日志里报出来
% 并指出替代写法。都能在 \pkg{caption} 的 \cs{captionsetup}\oarg{sub} 里找到对应：
%
% \begin{longtable}{ll}
% \toprule
% 旧写法 & 新写法 \\\midrule
% \endfirsthead
% \toprule
% 旧写法 & 新写法 \\\midrule
% \endhead
% \cs{subcaphangtrue}        & \cs{captionsetup}[sub]\{format=hang\} \\
% \cs{subcapcentertrue}      & \cs{captionsetup}[sub]\{justification=centering\} \\
% \cs{subcapcenterlasttrue}  & \cs{captionsetup}[sub]\{justification=centerlast\} \\
% \cs{subcapnoonelinetrue}   & \cs{captionsetup}[sub]\{singlelinecheck=false\} \\
% \cs{subcapraggedrighttrue} & \cs{captionsetup}[sub]\{justification=raggedright\} \\
% \cs{listsubcaptions}       & \cs{captionsetup}[sub]\{list=true\}（不自动转） \\\bottomrule
% \end{longtable}
%
% 把 \texttt{true} 换成 \texttt{false} 就是关掉，转成对应键的默认值。要注意
% \option{justification} 只有一个值，\cs{subcapcenter}、\cs{subcapcenterlast}、
% \cs{subcapraggedright} 三个开关都写它，后设的会盖掉先设的；\pkg{subfigure} 那边
% 是三个独立开关。
%
% \cs{listsubcaptions} 不自动转：它原先是在指定位置把攒下来的子图题一次性写进插图
% 索引，而 \option{list} 是每条子图题排出来时就写，等调用到 \cs{listsubcaptions}
% 已经晚了。用到只会提示，请自己在导言区写
% \cs{captionsetup}\oarg{sub}\marg{list=true}。
%
% \cs{subfigtopskip} 这一批长度仍然有效。
% 表格示例，表格中的字体是可以自行调整的。
% \begin{latex}
% \begin{table}[htbp]
% \caption[table1]{符合研究生院绘图规范的表格}[Table in agreement of the standard from graduate school]
% \vspace{0.5em}\centering\wuhao
% \begin{tabular}{ccccc}
% \toprule
% $D$(in) & $P_u$(lbs) & $u_u$(in) & $\beta$ & $G_f$(psi.in)\\
% \midrule
%  5 & 269.8 & 0.000674 & 1.79 & 0.04089\\
% 10 & 421.0 & 0.001035 & 3.59 & 0.04089\\
% 20 & 640.2 & 0.001565 & 7.18 & 0.04089\\
% \bottomrule
% \end{tabular}
% \end{table}
% \end{latex}
% 因为长表格不是浮动体，不会自动调整位置、也不会自动调整字体大小，一切都要手动设
% 置。特别繁琐。
% \begin{latex}
% \ltfontsize{\wuhao[1.667]} %设置表格内字体行间距
% \wuhao[1.667]\begin{longtable}{ccc} % 注意此处设置的是表格线距离
% \longbionenumcaption{}{{\wuhao 中国省级行政单位一览 %此处要添加字体设置
% }\label{table2}}{Table$\!$}{}{{\wuhao Overview of the provincial administrative
% unit of China}}{-0.5em}{3.15bp}\\ %注意后两个参数分别是中英标题间距、标题和表格的间距。
% %\caption{\wuhao 中国省级行政单位一览}\\[1em] %注意此处是标题和表格间距，这行
% %是单语标题
% \toprule 名称 & 简称 & 省会或首府  \\ \midrule
% \endfirsthead
% \multicolumn{3}{r}{表~\thetable（续表）}\vspace{0.5em}\\
% \toprule 名称 & 简称 & 省会或首府  \\ \midrule
% \endhead
% \bottomrule
% \endfoot
% 北京市 & 京 & 北京\\
% 天津市 & 津 & 天津\\
% 河北省 & 冀 & 石家庄市\\
% 山西省 & 晋 & 太原市\\
% 内蒙古自治区 & 蒙 & 呼和浩特市\\
% 辽宁省 & 辽 & 沈阳市\\
% 吉林省 & 吉 & 长春市\\
% 黑龙江省 & 黑 & 哈尔滨市\\
% 上海市 & 沪/申 & 上海\\
% 江苏省 & 苏 & 南京市\\
% 浙江省 & 浙 & 杭州市\\
% 安徽省 & 皖 & 合肥市\\
% 福建省 & 闽 & 福州市\\
% 江西省 & 赣 & 南昌市\\
% 山东省 & 鲁 & 济南市\\
% 河南省 & 豫 & 郑州市\\
% 湖北省 & 鄂 & 武汉市\\
% 湖南省 & 湘 & 长沙市\\
% 广东省 & 粤 & 广州市\\
% 广西壮族自治区 & 桂 & 南宁市\\
% 海南省 & 琼 & 海口市\\
% 重庆市 & 渝 & 重庆\\
% 四川省 & 川/蜀 & 成都市\\
% 贵州省 & 黔/贵 & 贵阳市\\
% 云南省 & 云/滇 & 昆明市\\
% 西藏自治区 & 藏 & 拉萨市\\
% 陕西省 & 陕/秦 & 西安市\\
% 甘肃省 & 甘/陇 & 兰州市\\
% 青海省 & 青 & 西宁市\\
% 宁夏回族自治区 & 宁 & 银川市\\
% 新疆维吾尔自治区 & 新 & 乌鲁木齐市\\
% 香港特别行政区 & 港 & 香港\\
% 澳门特别行政区 & 澳 & 澳门\\
% 台湾省 & 台 & 台北市\\
% \end{longtable}\normalsize %注意这里要恢复正常字体
% \end{latex}
% \subsubsection{公式}
% 公式不做介绍，与正常用法一致。
% \subsubsection{数学环境}
% \label{sec:math}
% \hithesis\ 定义了常用的数学环境：
%
% \begin{center}
% \begin{tabular}{*{7}{l}}\toprule
%   axiom & theorem & definition & proposition & lemma & conjecture &\\
%   公理 & 定理 & 定义 & 命题 & 引理 & 猜想 &\\\midrule
%   proof & corollary & example & exercise & assumption & remark & problem \\
%   证明 & 推论 & 例子& 练习 & 假设 & 注释 & 问题\\\bottomrule
% \end{tabular}
% \end{center}
%
% 比如：
% \begin{latex}
% \begin{definition}
%   道千乘之国，敬事而信，节用而爱人，使民以时。
% \end{definition}
% \end{latex}
% 产生（自动编号）：
% \medskip
%
% \noindent\framebox[\linewidth][l]{{\heiti 定义~1.1~~~} % {道千乘之国，敬事而信，节用而爱人，使民以时。}}
%
% \smallskip
% 列举出来的数学环境毕竟是有限的，如果想用\emph{胡说}这样的数学环境，那么可以定义：
% \begin{latex}
% \newtheorem{nonsense}{胡说}[chapter]
% \end{latex}
%
% 然后这样使用：
% \begin{latex}
% \begin{nonsense}
%   契丹武士要来中原夺武林秘笈。—— 慕容博
% \end{nonsense}
% \end{latex}
% 产生（自动编号）：
%
% \medskip
% \noindent\framebox[\linewidth][l]{{\heiti 胡说~1.1~~~} % {契丹武士要来中原夺武林秘笈。—— 慕容博}}
% \paragraph{调整定理的样式}
% \changes{v3.2a}{0000/00/00}{定理环境改用 \pkg{thmtools}，加 \option{theorem} 键族
%   与 \pkg{ntheorem} 命令的兼容层（\issue[183]）}
% 字体、间距、证毕符号这些写在 \option{theorem} 族里：
% \begin{latex}
% \hitsetup{
%   theorem = {
%     body-font  = {\songti\rmfamily},
%     head-font  = {\heiti\rmfamily},
%     qed        = {$\square$},
%     head-punct = {},
%   },
% }
% \end{latex}
%
% 键名与 \pkg{thmtools} 的 \cs{declaretheoremstyle} 一致，只是改成了连字符分词：
%
% \begin{longtable}{ll}
% \toprule
% 键 & 管什么 \\\midrule
% \endfirsthead
% \toprule
% 键 & 管什么 \\\midrule
% \endhead
% \option{body-font}   & 定理正文字体 \\
% \option{head-font}   & 定理标题字体 \\
% \option{note-font}   & 注记字体，留空则跟随标题 \\
% \option{note-braces} & 注记两侧的括号，两个参数 \\
% \option{head-punct}  & 标题末尾的标点 \\
% \option{head-space}  & 标题与正文之间的间距 \\
% \option{head-indent} & 标题的缩进 \\
% \option{qed}         & 证毕符号 \\
% \option{space-above} & 定理上方留白 \\
% \option{space-below} & 定理下方留白 \\
% \option{within}      & 编号跟随哪一级计数器 \\\bottomrule
% \end{longtable}
%
% 设置只对之后声明的环境生效，所以要写在 \cs{newtheorem} 之前。
%
% \paragraph{\pkg{ntheorem} 的写法}
% v3.2a 之前定理环境由 \pkg{ntheorem} 实现。\pkg{ntheorem} 的 \cs{@endtrivlist}
% 与 \pkg{enumitem} 冲突，\env{proof} 结尾的证毕符号会丢（\issue[292]），定理块
% 末尾还会多出一段空白（\issue[294]），所以改用了 \pkg{amsthm} 与 \pkg{thmtools}。
%
% 下面这些 \pkg{ntheorem} 专有命令仍然可用，编译时会在日志里提示新写法，
% v4.1（2027~年主要版本）移除：
%
% \begin{longtable}{ll}
% \toprule
% 旧命令 & 新写法 \\\midrule
% \endfirsthead
% \toprule
% 旧命令 & 新写法 \\\midrule
% \endhead
% \cs{theorembodyfont}   & \option{theorem} = \{ \option{body-font} \} \\
% \cs{theoremheaderfont} & \option{theorem} = \{ \option{head-font} \} \\
% \cs{theoremsymbol}     & \option{theorem} = \{ \option{qed} \} \\
% \cs{theoremseparator}  & \option{theorem} = \{ \option{head-punct} \} \\
% \cs{thref}             & \cs{autoref} \\
% \cs{listtheorems}      & \cs{listoftheorems} \\
% \cs{theoremlisttype}   & \cs{listoftheorems} 的可选参数 \\
% \cs{theoremprework}    & \cs{declaretheorem} 的钩子 \\
% \cs{theorempostwork}   & \cs{declaretheorem} 的钩子 \\
% \cs{theoremclass}      & \cs{declaretheorem} 的 \option{style} 键 \\
% \cs{theoremnumbering}  & \cs{newtheorem} 之后重定义 \cs{the}\meta{环境名} \\\bottomrule
% \end{longtable}
%
% \cs{theoremstyle} 现在是 \pkg{amsthm} 的命令，取值也换了一套，
% 原先 \pkg{ntheorem} 的 \texttt{break}、\texttt{change} 这些样式名不再认。
% 长度寄存器 \cs{theorempreskipamount} 与 \cs{theorempostskipamount} 也没有了，
% 请改用 \option{space-above} 与 \option{space-below}。
%
% \DescribeMacro{\quan}
% \DescribeMacro{\Quan}
% \DescribeMacro{\quansplit}
% \paragraph{带圈数字}
% \changes{v3.2a}{0000/00/00}{带圈数字优先用字体自带的字符，超出范围才画圈}
% 脚注号用的带圈数字由 \cs{quansplit}\marg{数字} 生成，也可以直接调用。当前宋体有
% 现成的带圈字符（\texttt{U+2460} 起，多数字体到 10 或 20）就用字符，字形比画出来的
% 规整；超出范围才用 \pkg{circledsteps} 画一个，并按字符版的尺寸缩放，混在一行里
% 看不出接缝。画的时候数字取当前宋体，与正文一致。
%
% \cs{quan}\oarg{十位}\marg{个位} 与 \cs{Quan}\marg{数字} 是同一个东西的两种写法，
% 压缩比例自动算，不必自己分位数。
%
% \subsubsection{算法}
% 我工算法不在规范中要求且一千个评审老师有一千个算法格式喜好。详见
% \href{https://github.com/PlutoThesis/PlutoThesis#%E6%B2%A1%E6%9C%89%E6%98%8E%E7%A1%AE%E8%A6%81%E6%B1%82%E7%9A%84%E6%A0%BC%E5%BC%8F}{PlutoThesis}
% 中的各个实验室算法喜好举例。在此多说无益。
% \subsubsection{引用参考文献}
% \DescribeMacro{\inlinecite}
% 学校要求的参考文献引用有两种模式：（1）上标模式。比如``同样的工作有很
% 多$^{[1,2]}$\ldots''。（2）正文模式。比如``文[3] 中详细说明了\ldots''。其中上标
% 模式使用远比正文模式频繁，所以为了符合使用习惯，上标模式仍然用常规
% 的 \cs{cite}\marg{key}，而 \cs{inlinecite}\marg{key} 则用来生成正文模式。
%
% 关于参考文献模板推荐使用 \BibTeX，关于中文参考文献需要额外增加一个 Entry:
% \texttt{language}，将其设置为 \texttt{zh} 用来指示此参考文献为中文，以
% 便 \file{hithesis.bst} 处理。如：
% \begin{latex}
% @INPROCEEDINGS{cnproceed,
%   author    = {王重阳 and 黄药师 and 欧阳峰 and 洪七公 and 段皇帝},
%   title     = {武林高手从入门到精通},
%   booktitle = {第~$N$~次华山论剑},
%   year      = 2006,
%   address   = {西安, 中国},
%   language      = "zh",
% }
%
% @ARTICLE{cnarticle,
%   AUTHOR  = "贾宝玉 and 林黛玉 and 薛宝钗 and 贾探春",
%   TITLE   = "论刘姥姥食量大如牛之现实意义",
%   JOURNAL = "红楼梦杂谈",
%   PAGES   = "260--266",
%   VOLUME  = "224",
%   YEAR    = "1800",
%   LANGUAGE    = "zh",
% }
% \end{latex}
%
% 注意如果不需要引用参考文献，请删除 \file{thesis.tex} 中 \cs{bibliography} 开头的两行，
% 以避免可能的编译错误。
%
% \subsubsection{列表环境}
% \DescribeEnv{itemize}
% \DescribeEnv{enumerate}
% \DescribeEnv{description}
% 为了适合中文习惯，模板将这三个常用的列表环境用 \pkg{enumitem} 进行了纵向间距压
% 缩。一方面清除了多余空间，另一方面用户可以自己指定列表环境的样式（如标签符号，
% 缩进等）。细节请参看 \pkg{enumitem} 文档，此处不再赘述。
% \subsection{后文}
%
% \subsubsection{结论}
% \DescribeEnv{conclusions}
% 结论之后为后文内容。
% \begin{latex}
% \begin{conclusions}
%
% 学位论文的结论作为论文正文的最后一章单独排写，但不加章标题序号。
%
% 结论应是作者在学位论文研究过程中所取得的创新性成果的概要总结，不能与摘要混为一
% 谈。博士学位论文结论应包括论文的主要结果、创新点、展望三部分，在结论中应概括论
% 文的核心观点，明确、客观地指出本研究内容的创新性成果（含新见解、新观点、方法创
% 新、技术创新、理论创新），并指出今后进一步在本研究方向进行研究工作的展望与设想
% 。对所取得的创新性成果应注意从定性和定量两方面给出科学、准确的评价，分（1）、
% （2）、（3）…条列出，宜用“提出了”、“建立了”等词叙述。
%
% \end{conclusions}
% \end{latex}
%
% \subsubsection{参考文献}
% 在后文中的参考文献是自动生成的，不需要用户干预，具体命令在 \file{thesis.tex} 中有
% 示例。
%
% \subsubsection{附录}
% \DescribeEnv{appendix}
% 所有的附录都插到这里来。因为附录会更改默认的 chapter 属性，而后面的{\heiti 个人简
%   历}又需要恢复，所以实现为环境可以保证全局的属性不受影响。
% \begin{latex}
% \begin{appendix}
% % -*-coding: utf-8 -*-
%%%%%%%%%%%%%%%%%%%%%%%%%%%%%%%%%%%%%%%%%%%%%%%%%%%%%%%%%
\chapter{带章节的附录}[Full Appendix]%
完整的附录内容，包含章节，公式，图表等

%%%%%%%%%%%%%%%%%%%%%%%%%%%%%%%%%%%%%%%%%%%%%%%%%%%%%%%%%
\section{附录节的内容}[Section in Appendix]
这是附录的节的内容

附录中图的示例：
\begin{figure}[htbp]
\centering
\includegraphics[width = 0.4\textwidth]{golfer}
%\bicaption[golfer5]{}{\xiaosi[0]打高尔夫球的人}{Fig.$\!$}{The person playing golf}\vspace{-1em}
\caption{\xiaosi[0]打高尔夫球的人}
\end{figure}

附录中公式的示例：
\begin{align}
a & = b \times c \\
E & = m c^2
\label{eq}
\end{align}

\chapter{这个星球上最好的免费Linux软件列表}[List of the Best Linux Software in our Planet]
\section*{系统}

\href{http://fvwm.org/}{FVWM 星球最强大的窗口管理器}——推荐

\section*{其他}

\href{https://github.com/goldendict/goldendict}{goldendict 星球最强大的桌面字典}——推荐

\href{http://www.mutt.org/}{mutt 星球最强大的邮件客户端}——推荐

% \end{appendix}
% \end{latex}
%
% \subsubsection{所发表文章}
% \DescribeEnv{publication}
% 虽然在\PGR\UGR\ 中都没有明确规定此处的格式，但按照旧模板PlutoThesis，此处格式
% 非常复杂。此处仍然使用旧模板中的设置方法。
% \begin{latex}
% \begin{publication}
% \noindent\textbf{（一）发表的学术论文}
% \begin{publist}
% \item	XXX，XXX. Static Oxidation Model of Al-Mg/C Dissipation Thermal Protection Materials[J]. Rare Metal Materials and Engineering, 2010, 39(Suppl. 1): 520-524.（SCI~收录，IDS号为~669JS，IF=0.16）
% \item XXX，XXX. 精密超声振动切削单晶铜的计算机仿真研究[J]. 系统仿真学报，2007，19（4）：738-741，753.（EI~收录号：20071310514841）
% \item XXX，XXX. 局部多孔质气体静压轴向轴承静态特性的数值求解[J]. 摩擦学学报，2007（1）：68-72.（EI~收录号：20071510544816）
% \item XXX，XXX. 硬脆光学晶体材料超精密切削理论研究综述[J]. 机械工程学报，2003，39（8）：15-22.（EI~收录号：2004088028875）
% \item XXX，XXX. 基于遗传算法的超精密切削加工表面粗糙度预测模型的参数辨识以及切削参数优化[J]. 机械工程学报，2005，41（11）：158-162.（EI~收录号：2006039650087）
% \item XXX，XXX. Discrete Sliding Mode Cintrok with Fuzzy Adaptive Reaching Law on 6-PEES Parallel Robot[C]. Intelligent System Design and Applications, Jinan, 2006: 649-652.（EI~收录号：20073210746529）
% \end{publist}
%
% \noindent\textbf{（二）申请及已获得的专利（无专利时此项不必列出）}
% \begin{publist}
% \item XXX，XXX. 一种温热外敷药制备方案：中国，88105607.3[P]. 1989-07-26.
% \end{publist}
%
% \noindent\textbf{（三）参与的科研项目及获奖情况}
% \begin{publist}
% \item	XXX，XXX. XX~气体静压轴承技术研究, XX~省自然科学基金项目.课题编号：XXXX.
% \item XXX，XXX. XX~静载下预应力混凝土房屋结构设计统一理论. 黑江省科学技术二等奖, 2007.
% \end{publist}
% %\vfill
% %\hangafter=1\hangindent=2em\noindent
% %\setlength{\parindent}{2em}
% \end{publication}
% \end{latex}
%
% \subsubsection{索引}
% \DescribeEnv{ceindex}
% 我工要求中英文双语索引。后文中的自动索引实际上不需要用户干预。
%\begin{latex}
% \begin{ceindex}
%   %如果想要手动加索引，注释掉以下这一样，用wordlist环境
% \printsubindex*
% \end{ceindex}
%\end{latex}
% 手工添加索引的方法不推荐，模板中将去除该功能。
% \subsubsection{授权}
% \DescribeMacro{\authorization}
% 授权页中的签名和日期是需要手写，不需要人工干预。具体示例在 \file{thesis.tex} 中。
%\begin{latex}
% \authorization %授权
% %\authorization[scan.pdf] %添加扫描页的命令，与上互斥
%\end{latex}
%
% \subsubsection{致谢声明}
% \DescribeEnv{acknowledgement}
% 把致谢做成一个环境更好一些，直接往里面写感谢的话就可以啦！
%
% \begin{latex}
% \begin{acknowledgement}
%   …
%   感谢\hit\LaTeX\ 论文模板\hithesis\ !
% \end{acknowledgement}
% \end{latex}
%
%
% \subsubsection{简历}
% \DescribeEnv{resume}
% 个人简历。
% 实际上，致谢和个人简历是自由发挥的地区，字体，文体，格式，内容，完全自己决定。
% \begin{latex}
% \begin{resume}
% XXXX~年~XX~月~XX~日出生于~XXXX。
%
% XXXX~年~XX~月考入~XX~大学~XX~院（系）XX~专业，XXXX~年~XX~月本科毕业并获得~XX~学学士学位。
%
% XXXX~年~XX~月------XXXX~年~XX~月在~XX~大学~XX~院（系）XX~学科学习并获得~XX~学硕士学位。
%
% XXXX~年~XX~月------XXXX~年~XX~月在~XX~大学~XX~院（系）XX~学科学习并获得~XX~学博士学位。
%
% 获奖情况：如获三好学生、优秀团干部、X~奖学金等（不含科研学术获奖）。
%
% 工作经历：
% \end{resume}
% \end{latex}
%
% \subsection{改写固定词汇}
% \label{sec:const}
% 封面标签、章节名、声明正文这类固定词汇由模板给出默认值。各院系的要求不尽相同，
% 需要改写时在导言区用 \option{const} 族设置：
% \begin{latex}
% \hitsetup{
%   const = {
%     zh = { supervisor-field-label = {指导教师}, keywords-separator = {；} },
%   },
% }
% \end{latex}
% 中英文成对的词汇同样收在 \option{zh} 与 \option{en} 分组里，写法见
% 第~\ref{sec:setup-language}~节。带后缀的本名与内部命令一一对应：
% \option{supervisor-field-label-zh} 对应 \cs{hit@supervisor@field@label@zh}，
% 名字里的连字符对应内部命令里的 \texttt{@}。
%
% v3.2a 之前这些词汇放在 \file{hithesisbook.cfg} 与 \file{hithesisart.cfg} 里，
% 需要改写就得改那两个文件。现在两个文件已经并入类文件，不再单独生成。
%
%
% \subsection{捐助}
% \changes{v1.0.01}{2017/08/27}{添加了捐助、矢量化本科论文模板的图片logo}
% 各位刀客和大侠如用的嗨，要解囊相助，请支付宝赞助
% 参照图~\ref{zfb}~中提示操作（二维码被矢量化后之后去
% 除了头像等冗余无用的部分～）。
% \begin{figure}[h]
% \centering\includegraphics[width=0.5\textwidth]{zfb.eps}
% \caption{支付宝}
% \label{zfb}
% \end{figure}
%
